\documentclass[11pt]{article}
\usepackage{amsmath,amssymb,amsthm}
\usepackage[margin=1in]{geometry}

\newtheorem*{theorem}{Theorem}

\begin{document}

\begin{theorem}
Let \(n\) be a positive integer. The probability that a uniform random function
\(f:[n]\to[n]\) has a unique cyclic vertex in its associated directed graph \(G_f\)
is \(1/n\).
\end{theorem}

\begin{proof}
Let
\[
\mathcal F_n=\{f:[n]\to[n]\}
\]
and let
\[
A_n=\{f\in\mathcal F_n: G_f\text{ has a unique cyclic vertex}\}.
\]
The map \(f\mapsto(f(1),\ldots,f(n))\) is a bijection from \(\mathcal F_n\) to
\([n]^n\), so \(|\mathcal F_n|=n^n\). By the uniform finite probability
convention,
\[
\Pr(A_n)=\frac{|A_n|}{|\mathcal F_n|}=\frac{|A_n|}{n^n}.
\]
It remains to count \(|A_n|\).

First record a finite-orbit fact. For \(f:[n]\to[n]\), define \(f^0(x)=x\) and
\(f^{k+1}(x)=f(f^k(x))\). For any \(x\in[n]\), the \(n+1\) vertices
\[
f^0(x),f^1(x),\ldots,f^n(x)
\]
lie in the \(n\)-element set \([n]\), so two are equal. If \(b>0\) is the least
index for which \(f^a(x)=f^b(x)\) for some \(a<b\), then
\[
f^a(x),f^{a+1}(x),\ldots,f^{b-1}(x)
\]
are pairwise distinct and form a directed cycle in \(G_f\). Hence every forward
orbit eventually enters a directed cycle. If \(G_f\) has a unique cyclic vertex
\(r\), then the cycle through \(r\) has only the vertex \(r\), so \(f(r)=r\).
Moreover, the eventual cycle reached by any vertex must be the loop at \(r\);
therefore every vertex reaches \(r\) after finitely many iterations of \(f\).

\medskip
\noindent\textbf{[KEY STEP].} Fix \(r\in[n]\), and set
\[
\mathcal A_{n,r}=\{f:[n]\to[n]\colon r\text{ is the unique cyclic vertex of }G_f\}.
\]
We prove that \(|\mathcal A_{n,r}|=n^{n-2}\) for \(n\ge2\), and that
\(|\mathcal A_{1,1}|=1\). If \(n=1\), the only function \([1]\to[1]\) has the
single loop at \(1\), so \(|\mathcal A_{1,1}|=1\).

Assume \(n\ge2\). We construct a bijection
\[
\Phi:\mathcal A_{n,r}\to [n]^{n-2}.
\]
For \(R\subseteq[n]\) containing \(r\), call \(a\in R\setminus\{r\}\) an
\(R\)-leaf if no vertex \(y\in R\setminus\{a\}\) has \(f(y)=a\).

Given \(f\in\mathcal A_{n,r}\), start with \(R_1=[n]\). For
\(k=1,\ldots,n-2\), choose \(a_k\) to be the least \(R_k\)-leaf in
\(R_k\setminus\{r\}\), set \(w_k=f(a_k)\), and put
\(R_{k+1}=R_k\setminus\{a_k\}\). This is well-defined. Indeed, suppose at the
beginning of a step that \(r\in R_k\), that \(f(R_k)\subseteq R_k\), and that
every vertex of \(R_k\) reaches \(r\). If no non-root \(R_k\)-leaf existed, then
for every \(u\in R_k\setminus\{r\}\) there would be a vertex
\(p(u)\in R_k\setminus\{u\}\) with \(f(p(u))=u\). Since \(f(r)=r\), this
\(p(u)\) lies in \(R_k\setminus\{r\}\). Iterating \(p\) in the finite set
\(R_k\setminus\{r\}\) would repeat and produce a directed cycle for \(f\)
entirely outside \(r\), contradicting the fact that every remaining vertex
reaches \(r\). Thus a least non-root \(R_k\)-leaf exists.

For such a leaf \(a_k\), one has \(f(a_k)\ne a_k\), since otherwise
\(a_k\ne r\) would be cyclic. Thus no vertex of \(R_k\) maps to \(a_k\), and
deleting \(a_k\) preserves \(f(R_{k+1})\subseteq R_{k+1}\). The remaining
vertices still reach \(r\), because a path from a remaining vertex to \(r\)
cannot pass through \(a_k\) without some remaining vertex mapping to \(a_k\).
The induction continues until two vertices remain, \(R_{n-1}=\{r,b\}\), and then
closure, \(f(r)=r\), and reachability force \(f(b)=r\). Define
\[
\Phi(f)=(w_1,\ldots,w_{n-2}).
\]

Conversely, given \(w=(w_1,\ldots,w_{n-2})\in[n]^{n-2}\), start with
\(R_1=[n]\). For \(k=1,\ldots,n-2\), choose \(a_k\) to be the least element of
\(R_k\setminus\{r\}\) that does not occur among
\[
w_k,w_{k+1},\ldots,w_{n-2}.
\]
Such an element exists because \(R_k\setminus\{r\}\) has \(n-k\) elements and
the suffix has length \(n-k-1\). Define \(f(a_k)=w_k\) and set
\(R_{k+1}=R_k\setminus\{a_k\}\). The assignment is legitimate: \(w_k\ne a_k\),
and \(w_k\) was not removed earlier, for if \(w_k=a_j\) with \(j<k\), then
\(a_j\) would have appeared in the suffix \(w_j,\ldots,w_{n-2}\), contrary to
the choice of \(a_j\). After the \(n-2\) steps, write the remaining set as
\(\{r,b\}\), and define \(f(b)=r\) and \(f(r)=r\). This defines a function
\(f:[n]\to[n]\). If \(a_k\) has index \(k\), \(b\) has index \(n-1\), and
\(r\) has index \(n\), then \(f(v)\) has strictly larger index than \(v\) for
every \(v\ne r\). Hence every non-root vertex reaches \(r\), while \(r\) is
fixed, so the only directed cycle is the loop at \(r\). Thus the decoding gives
a map \(\Psi:[n]^{n-2}\to\mathcal A_{n,r}\).

The maps \(\Phi\) and \(\Psi\) are inverse maps. Let \(f\in\mathcal A_{n,r}\)
and let \(\Phi(f)=w\). At encoding step \(k\), for \(y\in R_k\setminus\{r\}\),
\[
y\text{ is an }R_k\text{-leaf}
\quad\Longleftrightarrow\quad
y\text{ does not occur among }w_k,\ldots,w_{n-2}.
\]
If \(y=w_j=f(a_j)\) for some \(j\ge k\), then the still-present vertex
\(a_j\ne y\) maps to \(y\), so \(y\) is not a leaf. Conversely, if \(y\) is not
a leaf, choose \(z\in R_k\setminus\{y\}\) with \(f(z)=y\). Since \(y\ne r\),
the vertex \(z\) is neither \(r\) nor the final vertex \(b\), because
\(f(r)=f(b)=r\). Thus \(z=a_j\) for some \(j\ge k\), and \(w_j=y\). Therefore
decoding \(\Phi(f)\) chooses the same least vertex at every step and reconstructs
\(f\), so \(\Psi(\Phi(f))=f\).

Conversely, let \(w\in[n]^{n-2}\) and let \(\Psi(w)=f\). At decoding step \(k\),
the chosen \(a_k\) is an \(R_k\)-leaf, since it does not appear in the remaining
suffix and neither \(b\) nor \(r\) maps to it. If \(y<a_k\) is a non-root
remaining vertex, then by the choice of \(a_k\), the vertex \(y\) occurs in the
remaining suffix, so some still-present vertex maps to \(y\). Hence \(a_k\) is
the least \(R_k\)-leaf, and the encoding records \(f(a_k)=w_k\) at each step.
Thus \(\Phi(\Psi(w))=w\).

Therefore \(\Phi\) is a bijection from \(\mathcal A_{n,r}\) to \([n]^{n-2}\).
Since there are \(n^{n-2}\) words of length \(n-2\) over \([n]\), we have
\[
|\mathcal A_{n,r}|=n^{n-2}\qquad(n\ge2).
\]

Now for each \(r\in[n]\), let
\[
A_{n,r}=\{f\in\mathcal F_n: r\text{ is the unique cyclic vertex of }G_f\}.
\]
Then
\[
A_n=\bigsqcup_{r\in[n]}A_{n,r},
\]
because every function with a unique cyclic vertex has exactly one such vertex
\(r\in[n]\), and a unique vertex cannot be two distinct elements. If \(n=1\),
then \(|A_n|=1=1^{1-1}\). If \(n\ge2\), then each fixed-root class has size
\(n^{n-2}\), so
\[
|A_n|=\sum_{r\in[n]}|A_{n,r}|=n\cdot n^{n-2}=n^{n-1}.
\]
Thus \(|A_n|=n^{n-1}\) for every positive integer \(n\). Finally,
\[
\Pr(A_n)=\frac{|A_n|}{n^n}
=\frac{n^{n-1}}{n^n}
=\frac1n.
\]
This proves the theorem.
\end{proof}

\end{document}
